\documentclass[titlepage, a4paper,12pt, twoside]{article}

%TC:macro \cite [option:ignore]
%TC:macro \citep [option:ignore]
%TC:macro \citet [option:ignore]
%TC:macro \ref [option:ignore]
%TC:macro \eqref [option:ignore]
%TC:macro \label [option:ignore]

%TC:envir abstract 0 0
%TC:envir acknowledgements 0 0
%TC:envir acknowledgments 0 0
%TC:envir appendix 0 0
%TC:envir figure 0 0
%TC:envir table 0 0
%TC:envir equation 0 0
%TC:envir tabular 0 0

%TC:macro \caption [option:ignore]
%TC:macro \maketitle [option:ignore]
%TC:macro \tableofcontents [option:ignore]
%TC:macro \listoffigures [option:ignore]
%TC:macro \listoftables [option:ignore]

%TC:macro \begin [option:ignore]
%TC:macro \end [option:ignore]

%TC:break bibliography 0 0
%TC:break thebibliography 0 0
%TC:break references 0 0

\usepackage{lipsum}

\usepackage{csquotes}
\usepackage{svg}
\usepackage[english]{babel}
\usepackage{rotating}
\usepackage{url}
\usepackage{amsmath}
\usepackage{amssymb}
\usepackage{amsfonts}
\usepackage{braket}
\usepackage{todonotes}

\usepackage{graphicx}
\graphicspath{{figs/}}

\usepackage{cases}
\usepackage{float}
\usepackage{hyperref}
% todo figures
\usepackage{caption}

\usepackage[margin=2cm, bindingoffset=0.5cm, top=4cm]{geometry}
\usepackage{subcaption}
\usepackage{fancyhdr}

\fancyhf{}
\setlength{\headheight}{15pt}
% --- Header Configuration ---
% E = Even page, O = Odd page
% L = Left, R = Right

% Even pages (Left side of the spread)
\fancyhead[LE]{\thepage}          % Page number: Outer
\fancyhead[RE]{\textit{Thomas Holland}}        % Author: Inner

% Odd pages (Right side of the spread)
\fancyhead[LO]{\textit{Methods for Estimating Ice Melt Driven Sea Level Change}}       % Short Title: Inner
\fancyhead[RO]{\thepage}          % Page number: Outer

\renewcommand{\headrulewidth}{0pt} % Removes the horizontal line

\usepackage[T1]{fontenc}


\usepackage[backend=bibtex, uniquename=full, maxnames=2, style=authoryear]{biblatex}
\addbibresource{zotero.bib}

\usepackage{booktabs}
\usepackage{tabularx}
\renewcommand{\tabularxcolumn}[1]{>{\raggedright\arraybackslash}p{#1}}

\let\oldparagraph\paragraph
\renewcommand{\paragraph}[1]{\oldparagraph{#1}\mbox{}\\}

\let\stdsection\section
\renewcommand
\section{\newpage\stdsection}


\newcommand{\sshc}{\Delta \mathrm{SSH}}
\newcommand{\slc}{\Delta \mathrm{SL}}
\newcommand{\gmslc}{\Delta \mathrm{GMSL}}
\newcommand{\gmsl}{\mathrm{GMSL}}
\newcommand{\altimetry}{\mathrm{ALT}}
\newcommand{\odtchange}{\Delta \mathrm{SSH}_{\mathrm{OD}}}

% gmsl estimate variants
\newcommand{\gmslctrue}{\Delta \mathrm{GMSL}^{\mathrm{true}}}
\newcommand{\gmslcbay}{\Delta \mathrm{GMSL}^{\mathrm{bay}}}
\newcommand{\gmslcalt}{\Delta \mathrm{GMSL}^{\mathrm{alt}}}

\newcommand{\covariance}[1]{\mathcal{C}_{#1}}
\newcommand{\thicknesschange}[1]{\Delta T_{#1}}
\newcommand{\masschange}[1]{\Delta M_{#1}}
\newcommand{\operator}[2]{\operatorname{{#1}}_{{#2}}}

% covariances


% mass changes


% operators
\newcommand{\operatorsshc}{\operator{F}{\sshc}}
\newcommand{\operatorslc}{\operator{F}{\slc}}
\newcommand{\operatorloadwater}{\operator{L}{w}}
\newcommand{\operatorloadice}{\operator{L}{i}}
\newcommand{\operatorloadfirn}{\operator{L}{f}}
\newcommand{\operatorgmsl}[1]{\operator{G}{#1}}
\newcommand{\operatoraltimetry}[1]{\operator{A}{#1}}

% point-sampling operators (altimetry subscript variants)
\newcommand{\operatorpointocean}{\operatoraltimetry{S}}
\newcommand{\operatorpointtg}{\operatoraltimetry{T}}
\newcommand{\operatorpointice}{\operatoraltimetry{I}}


\begin{document}

%TC:ignore

\begin{titlepage}
  \newgeometry{margin=2.8cm}

  \begin{center}
    \vspace*{0.5cm}

    \makebox[\textwidth][l]{\includegraphics[width=0.55\textwidth]{ucam_dep_of_earth_sciences_h_col.pdf}}

    \vspace{2.2cm}

    {\LARGE\bfseries
      Improving Methods for\\[0.15em]
      Estimating Modern Ice-Driven\\[0.15em]
      Sea Level Change\\[0.15em]
      using Bayesian Inference\par}

    \vspace{2.1cm}

    \includegraphics[width=3.6cm]{crest.pdf}

    \vspace{1.7cm}

    {\Large\bfseries Thomas Holland\par}

    \vspace{0.18cm}

    {\normalsize
      Jesus College
      \par}

    \vfill

    {\large
      Supervised by Dr David Al-Attar
      \par}

    \vfill

    {\normalsize
      Project submitted for degree of\\[0.2em]
      \textit{Natural Sciences MSci (Earth Sciences)}
      \par}

    \vfill

    {\large March 2026\par}

  \end{center}

  \restoregeometry
\end{titlepage}

\clearpage
\thispagestyle{empty}
\newgeometry{margin=1.2cm}

\begin{center}
  \vspace*{\fill}
  \setlength{\tabcolsep}{6pt}
  \renewcommand{\arraystretch}{0}
  \begin{tabular}{cc}
    \includegraphics[width=0.29\textwidth]{6-9_posterior_ice_thickness_antarctica.pdf}      &
    \includegraphics[width=0.29\textwidth]{6-11_posterior_ice_thickness_greenland.pdf}        \\
    \includegraphics[width=0.29\textwidth]{joint_precon_grace_ice_firn_gmsl_covariance.pdf} &
    \includegraphics[width=0.29\textwidth]{z_score_comparison_contour_a.pdf}
  \end{tabular}

  \vspace{0.3cm}

  \includegraphics[width=0.44\textwidth]{major_source_altimetry_errors_scalar.pdf}
  \vspace*{\fill}

  \begin{minipage}{0.9\textwidth}
    \centering
    \footnotesize
    This report is submitted for the degree of Natural Sciences MSci. (Earth Sciences).

    \vspace{6pt}

    {7489 words.}

    \vspace{6pt}

    All work and data found within this report can be found here:\\
    \url{https://github.com/0thomasholland/Part_III_Project}

    References to specific code notebooks are made throughout via links and references section.

    \vspace{6pt}

    I declare that the submitted work is my own, except where acknowledgement is given to the work of others or to work done in collaboration. Code authorship is my own and verifiable via the GitHub link above. The work within builds upon the work done by Heathcote and Al-Attar (2026), and so naturally some figures and method validation are as such similar. I used contextual autocomplete within Zed for programming and Claude (Anthropic) for broad assists with improving clarity and structure in writing.\\ All substantive ideas, analyses, and conclusions are my own.
  \end{minipage}
\end{center}

\restoregeometry


\begin{abstract}
Satellite altimetry is widely used to estimate global mean sea level change (GMSLC), but the standard approach treats sea surface height change (SSHC) as a proxy for sea level change (SLC), neglecting the gravitationally self-consistent redistribution of water that links ice mass loss to the sea surface. This report quantifies the systematic bias introduced by this proxy and develops a principled Bayesian alternative.

Error quantification for the three major ice sheets shows that the SSHC proxy introduces source-dependent biases: GMSLC is underestimated by approximately 5.8--8.2\% for Greenland melt and 0.5--2\% for East Antarctic melt, and overestimated by approximately 1--2\% for West Antarctic melt. Because the bias depends on where ice is melting, no scalar correction can account for it; correcting for the error requires knowledge of the very quantity being estimated.

In place of the proxy, we apply the infinite-dimensional Bayesian inversion framework of \textcite{heathcoteScaleableBayesianFramework2026}, implemented in \citetitle{al-attarPygeoinfPythonLibrary2025}, to recover ice thickness change directly from SSH altimetry. Priors are defined as Gaussian measures over Sobolev spaces on the sphere, encoding physically motivated beliefs about the spatial structure of mass change via a novel Richards curve construction conditioned on ice thickness. The posterior is obtained in closed form and propagated through the full sea level physics without discretisation artefacts. Across 261 synthetic twin experiments, 60.5\% of true GMSLC values fall within one posterior standard deviation, compared to 14.6\% under the proxy method, providing direct empirical evidence of well-calibrated uncertainty quantification.

The inversion is extended to a joint inference over ice thickness change, firn compaction, and ocean dynamic topography, combining SSH altimetry, ice altimetry, and GRACE gravity observations in a single composite model space. The joint forward operator is factorised to reduce computational cost by a factor of three without approximation. Knock-out tests demonstrate that each data source contributes non-redundant information, and the posterior successfully disentangles contributions that are entangled at the level of the observations. Sensitivity tests confirm robustness to prior mis-specification in the inversion results, though computational efficiency is sensitive to the prior covariance amplitude.

Taken together, the methods developed here offer three technical advances over both the SSHC proxy and existing discrete parameterisation approaches: principled uncertainty propagation through the full forward physics; spatially correlated posteriors whose structure is inherited from physically motivated priors rather than imposed by a discretisation grid; and joint inference over competing signals with formal uncertainty quantification on the trade-offs between them.
\end{abstract}

\pagestyle{empty}
\tableofcontents
\newpage
\pagestyle{fancy}
%TC:endignore

\section{Introduction}
\todo{this whole section}

% Glacial driven sea level change

This report investigates and provides an improved novel method to estimate modern sea level from ice sheet melt, based on sea level physics.

Sea Level (\textbf{SL}), also known as Relative Sea Level, is the difference in elevation between the sea surface and the solid earth surface. A change in this elevation difference is termed Sea Level Change (\textbf{SLC}).

Sea Surface Height (\textbf{SSH}), also known as Absolute Sea Level, is defined as the elevation of the sea surface above a reference ellipsoid. Ocean satellite altimetry measures Sea Surface Height, as it measures the height difference between itself on an orbital ellipsoid and the sea surface.

 However, the quantity of primary interest is SLC, as this is what is directly relevant to coastal communities and policy-makers.




% Importance

Understanding modern SLC is important for blah. Global Mean Sea Level (\textbf{GMSL}) is the areal average height of SL \autocite{hortonMappingSeaLevelChange2018}. GMSL is a particularly important quantity as it decreases the short timescale effects that affect SL, and allows for tracking of global trends. It is genererally accepted that GMSL change is to a large part influenced by ice sheet and glacial mass flux signals \autocite[e.g.][]{slangenEvaluatingModelSimulations2017}.


Smth regional SLC


% Idea

Ice melt driven SLC has an rotationally vairant impact on the sea surface.

This means that as well as being able to calcualte GMSL from SLC, it is possible to calcualte where the ice sheet change occured to drive the SLC.

Can be useful to understand where melt is occuring etc

% traditional method

take a surface average of SSH and apply correction (?) and call it SLC

% Problem

errors assocaited with traditional are poorly constrained and come from two main places.

first is there is the intrinsic

also, other assumptive corrections are needed to account for non-ice driven GMSL changes, such as.

also, investigation into density variations haven't been investigate


the errors cannot be accoutned for with a simple correction as they are depenednt on the ice source location

believed that the errors are sigificant and locationally variable


% Solution

By using a framework presented in \textcite{heathcoteScaleableBayesianFramework2026}, it is possible to explore modern ice driven sea level change using the full account of the physics related to the gravitationally consistent sea level changein a infinite-demiensional setting. This allows for novel quantification of the errors assocaited with traditional methods of GMSL


REF DA has shown that using a Bayseian inversion method with tide gauge data can lead to more

as includes physical aspects of the problem built into the problem

by including the physics of the rebound effect and other noise signals belived to increase accuracy

\subsection{Sea level physics in relation to ice sheet change}\label{sec:slp}

This will act as an introduction to how Sea Level is affected by ice sheet change.

SL defines the sea surface as an equipotential surface:

\begin{equation}
	\Phi(\mathbf{x} + \SL \hat{\mathbf{r}}) = \Phi_0 \label{eq:sea_level_equipotential}
\end{equation}

When this is perturbed by an ice thickness change, $\thicknesschangeice$, three things occur simultaneously: the gravitational potential itself changes, by $\phi$; the solid Earth surface, $\mathbf{x}$, moves radially, by $\mathrm{u}$; the value of the potential at the ocean surface shifts, from $\Phi_0$ to $\Phi_0 + \Phi_g$. This results in the SL changing, from $\SL$ to $\SL + \Delta \SL$, to maintain the equipotential condition. This can be captured within the following equation:

\begin{equation}
	(\Phi + \phi)(\mathbf{x} + \mathbf{u} + (\SL + \slc)\mathbf{\hat{r}}) = \Phi_0 + \Phi_g
\end{equation}

If we apply a first order Taylor expansion to the left-hand side, we get

\begin{equation}
  \Phi(\mathbf{x} + \SL \hat{\mathbf{r}}) + \phi + \nabla \Phi \cdot (\mathbf{u} + \slc \hat{\mathbf{r}}) = \Phi_0 + \Phi_g
\end{equation}


which can be simplfied: using \autoref{eq:sea_level_equipotential} we can cancel these terms from both sides; $\nabla\Phi$ represents the gravitational acceleration vector with its radial magnitude is $g$, therefore $\nabla\Phi \cdot \mathbf{\hat{r}} = g$; term $u \cdot \nabla\Phi$ represents the potential change due to the vertical motion of the solid surface:

\begin{equation}
	\phi + u \cdot \nabla\Phi + g \slc = \Phi_g
\end{equation}

Which we can rearrange to get the change in sea level:

\begin{equation}
	\slc = -\frac{1}{g}(u \cdot \nabla\Phi + \phi) + \frac{\Phi_g}{g} \label{eq:sl-equation}
\end{equation}

When discussing SSH, we need to account for the fact that SSH is the height of the sea surface above a reference ellipsoid, which is not an equipotential surface. This means that the change in SSH, $\sshc$, is related to the change in SL, $\slc$, by:

\begin{equation}
  \sshc = \slc + \frac{1}{g}(\mathbf{u} \cdot \nabla\Phi + \phi) \label{eq:sshc_slc_relation}
\end{equation}

where $\frac{1}{g}(\mathbf{u} \cdot \nabla\Phi)$ accounts for the vertical crustal motion and $\frac{1}{g}\phi$ accounts for the change in gravitational potential. Using \autoref{eq:sl-equation} this results in:

\begin{equation}
  \sshc = \frac{\Phi_g}{g}
\end{equation}



If we have an ice load, $\zeta$, which is due to the change in ice thickness, $\thicknesschangeice$ with density $\densityice$, then the ice load is:

\begin{equation}
	\zeta = \densityice (1-C) \thicknesschange{I}
\end{equation}

Where $1 - C$ accounts for the possibility of floating ice and C is the ocean function, which is 1 if the point is ocean and 0 if the point is land. This means that the ice load only applies to points where there is no floating ice, as floating ice does not contribute to the load on the solid earth.





\todo{fix this}
The sea surface, when ignoring shorter term dynamic effects, lies on an equipotential gravity surface. Ice mass change results in two effects to this equipotential surface, firstly the added water volume shifts the surface from one equipotential level to another, and secondariliy the change in mass results in

This can be captured within the following equation



Ocean function is
\begin{align}
	\mathrm{C}(\mathbf{x})= \begin{cases}
		1, &\quad\rho_w\mathrm{SL}(x)-\rho_i\mathrm{I}(x)>0 \\
		0, &\quad\rho_w\mathrm{SL}(x)-\rho_i\mathrm{I}(x)\le0
			\end{cases}
\end{align}

SLC can then be defined as:

\begin{align}
	\Delta SL = -\frac{\Delta M_\text{ice}}{\rho_w A}
\end{align}

Where:
\begin{align*}
	A &= \int_{\partial M}\mathrm{C}\, dS \\ &= r^2\int_{\mathbb{S}^2}\mathrm{C}\, dS
\end{align*}


For Glacial isostatic adjustment, relative sea level history, ice thickness history, and mantle viscosity \citetitle{royICE7G2018} \autocite{royICE7G2018} model is used see \textcite{royGlacialIsostaticAdjustment2015} and \textcite{al-attarPySLFPPythonSea2025} for more details.


\section{Background to mathematical methods}
The methods presented here rely on a mathematical framework implemented in the \citetitle{al-attarPygeoinfPythonLibrary2025} python package as described in \textcite{heathcoteScaleableBayesianFramework2026}. This framework represents continuous fields as functions upon a sphere; applies linear operators to these fields; and uses Gaussian measures to quantify uncertainty. The mathematical properties of this framework are critical to the development and implementation of the methods described in this report, and so a brief overview is given here to provide context for the subsequent sections.

For most readers, the framing given should be sufficient to understand the remainder of the paper. Full details of the mathematical framework surrounding the underlying python packages can be found in \textcite{heathcoteScaleableBayesianFramework2026}.

\subsection{Mathematical framework}\label{sec:maths_background}

For sea-level physics at the relevant lengthscales and timescales, it is physically appropriate to pose our problem on spaces where the realised fields are guaranteed to be spatially smooth, and where linear operators represent the physical mapping between these spaces.

To represent a continuous field over a sphere ($\mathbb{S}^2$, scaled to the dimensions of the Earth) we work within a Sobolev space, $H^s\left(\mathbb{S}^2\right)$, where the index $s$ controls the smoothness of admissible fields.

Linear operators (e.g. $A$) map between these spaces, with their adjoints mapping back to the original space (e.g. $A^*$). For example, if we have a linear operator $A: X \to Y$ that maps from a space $X$ to a space $Y$, then its adjoint $A^*: Y \to X$ maps back from $Y$ to $X$.

This framework accommodates both deterministic fields, which enter as fixed elements of $H^s\left(\mathbb{S}^2\right)$, and Gaussian measures, which are used to quantify our understanding of various parameters, as described later in \autoref{sec:bay}. These Gaussian measures are characterised by their mean, $m \in H^s\left(\mathbb{S}^2\right)$, and covariance operator, $\mathcal{C}: H^s\left(\mathbb{S}^2\right) \to H^s\left(\mathbb{S}^2\right)$, and are written as

\begin{equation}
\mu = \mathcal{N}\left(m, \mathcal{C}\right).
\end{equation}

There are suitable conditions that $\mathcal{C}$ must satisfy to be a valid covariance operator on an infinite-dimensional space, which are handled within \citetitle{al-attarPygeoinfPythonLibrary2025}. An example is the heat kernel covariance structure. Let $\Delta$ be the Laplace-Beltrami operator on the sphere, $\heatamp$ and $\heatlen$ parameters controlling the pointwise standard deviation and the lengthscale of spatial correlations respectively, and $\braket{\cdot,\cdot}$ the inner product of the space (a generalisation of the dot product \autocite{parkerGeophysicalInverseTheory1994}). The heat kernel covariance then takes the form

\begin{equation}
  \covariance{} = \frac{\heatamp^2}{\braket{e^{-\heatlen^2\Delta}\delta_x,\delta_x}} e^{-\heatlen^2 \Delta}\label{eq:heat_kernel_covariance}
\end{equation}

This covariance structure ensures the resulting Gaussian measure is well-defined and rotationally invariant, meaning the covariance between any two points depends only on their geodesic separation rather than their absolute position on the sphere. It is useful for geophysical problems, encoding the physical intuition of scaled point-wise variance and spatial correlation that decays with distance.

Linear operators acting on Gaussian measures produce a valid Gaussian pushforward measure on the target space:

\begin{align}
\mu_X &= \mathcal{N}\left(m, \mathcal{C}\right) \in X, \\
A\mu_X &= \mathcal{N}\left(Am, A\mathcal{C}A^*\right) \in Y.
\end{align}

That is, pushing $\mu_X$ forward through $A$ transforms the mean and covariance consistently; this closure under linear maps permits the posterior distribution to be computed in closed form, as shown in \autoref{sec:bay}.

It is often the case that we wish to extract scalar-valued quantities of interest, such as the Global Mean Sea Level (GMSL). These can be represented as continuous linear functionals $\ell: H^s\left(\mathbb{S}^2\right) \rightarrow \mathbb{R}$. By the Riesz representation theorem \autocite{rudinFunctionalAnalysis1991}, any such functional can be written as

\begin{equation}
\ell(h) = \braket{h, k}_{H^s}
\end{equation}

for a unique representer $k \in H^s\left(\mathbb{S}^2\right)$. This allows us to compute the distribution of GMSLC directly from the distribution of SLC.

\subsection{Bayesian inference framework}\label{sec:bay}

Recovering ice sheet change from altimetry is an inverse problem, which we address here using Bayesian inference. This implements Bayes' theorem as:

\begin{equation}
	p\left(\text{model} | \text{data}\right) = \dfrac{p\left(\text{data} | \text{model}\right)p\left(\text{model}\right)}{p\left(\text{data}\right)}
\end{equation}

where $p\left(\cdot\right)$ denotes probability, the data are the observations we collect or synthetically generate when developing the methods, and the model comprises the parameters we define across the model space. The likelihood, $p\left(\text{data} | \text{model}\right)$, is the probability of observing the data given a particular model. The prior, $p\left(\text{model}\right)$, is our initial belief about the model before seeing the data. The posterior, $p\left(\text{model} | \text{data}\right)$, is our updated belief after observing the data.

Bayesian inversion applies this framework to inverse problems. Rather than seeking a single best-fit solution, it characterises the entire probability distribution of unknown physical parameters from noisy data. Concretely, if we express the prior over the model space as a Gaussian measure with expectation $\bar{m}$ and covariance $\mathcal{C}_m$, the model space can be represented as

\begin{equation*}
	m \sim \mathcal{N}\left(\bar{m}, \mathcal{C}_m\right)
\end{equation*}

which we push forward using an operator $A$ that represents the physical relation between model and data spaces. Measurements carry random noise, modelled as a Gaussian vector $n$ with expectation $\bar{n}$ and covariance $\mathcal{C}_n$. The data vector $d$ then is

\begin{align*}
	d &= Am +n,\\
	\implies  d &\sim \mathcal{N}\left(A\bar{m} + \bar{n},\quad A\mathcal{C}_mA^* + \mathcal{C}_n\right).
\end{align*}

Applying the Gaussian inverse problem form of Bayes' theorem, first presented in \textcite{stuartInverseProblemsBayesian2010} and implemented in \textcite{heathcoteScaleableBayesianFramework2026}, the posterior is available in closed form as a Gaussian measure with mean and covariance

\begin{align}
	m_p \sim \mathcal{N}\left(\bar{m}+K\left(\dobs-A\bar{m}-\bar{n}\right),\quad \left(1-KA\right)\mathcal{C}_{m}\right)
\end{align}

where $K$ is the Kalman gain, defined as

\begin{equation}
	K = \mathcal{C}_{m}A^*\left(A\mathcal{C}_{m}A^*+\mathcal{C}_{n}\right)^{-1}.
\end{equation}

Evaluating these expressions requires inverting the normal operator $N = A\mathcal{C}_{m}A^* + \mathcal{C}_{n}$, which is performed iteratively via conjugate gradients. Each action of $N$ requires one forward and one adjoint application of $A$; the posterior expressions themselves involve no approximation.

Priors defined in the model space are critical to the inversion results, as they encode our beliefs about the system and significantly influence the posterior. Priors that encode physical knowledge of the system are called informative priors. An example is using spatial patterns of ice thickness change or constraints based on physical laws to restrict prior distributions. These informative priors can help to regularise the inversion and guide it towards physically plausible solutions, especially in cases where the data is sparse or very noisy.

Using informative priors, however, is not without its flaws. We must assess how informative priors affect the posterior distribution, as poorly chosen ones can introduce bias \autocite{vandeschootBayesianStatisticsModelling2021}. The distinction is nuanced. For our purposes: restrictions motivated by physical principles are valid; restrictions motivated by what we believe has generated the observed data are not. For example, restricting ice loss to the margins of ice sheets is physically motivated and so constitutes a valid informative prior. However, biasing these priors towards an expectation of ice loss based on prior knowledge of the observations is not, as it introduces a bias driven by what we expect the data to show.


\section{Methods}\label{sec:methods}
\subsection{Underlying methods}

The project builds upon and utilises REF pyslfp which itself is built upon pygeoinf \cite{al-attar2023}

Mathematical methods

Fields of either guassian fields or scalars

linear operators defined to map from one to another, with their adjoint (the Hermitian) mapping to the inverse

\begin{align}
  A: X &\rightarrow Y \text{ such that}\\
  y &= A(x) \\
  A^*: Y &\rightarrow X \text{ such that}\\
  x&= A^*(y)
\end{align}


\subsection{Error Quantification}

WHole ice sheets

major sources

major sources mixing

\subsection{Bayseian Inversion}

% Breif introduction

By applying the inversion framework outlined in \autoref{sec:bay} we can combined our understanding of SL physics and with prior distributions built to reflect our knowledge of the system to result in an updated posterior belief of the state of the model space based on observations made in the data space.



\subsubsection{Simple inversion}

\notebookcite{thomasholland06SimpleInversionipynb2026}

A simple case for this is to invert for ice thickness change, $\thicknesschange{i}$, with a single density, from SSH altimetery, $d$ where the forward problem takes the form

\begin{align}
	d &= \operatoraltimetry{\text{ocean}} \circ \operatorsshc \circ \operatorslc \circ \operatorloadice (\thicknesschange{i}) + n,\\
	d &= \operator{O}{simple}(\thicknesschange{i}) + n, \label{eq:simple_case}
\end{align}

where $\operator{O}{simple}$ is the composed forward operator for this simple case and $n$ represents the altimetry noise. The noise is taken to be an independent Gaussian measure in the dataspace with zero mean and covariance $\covariance{n}$, so

\begin{equation}
	n \sim \mathcal{N}(0, \covariance{n}).
\end{equation}

A true ice-thickness field, $\thicknesschange{i}^\dagger$,  is drawn as a single sample from the prior measure $\thicknesschange{i}$, providing a plausible but non-trivial test case that is, by construction, consistent with our prior assumptions. Synthetic altimetry observations, $d_{obs}$, are then generated by passing through the forward operator and adding the Gaussian noise:

\begin{equation}
	d_{obs} = \operator{O}{simple}(\thicknesschange{i}^\dagger) + n
\end{equation}

Using the Bayesian inversion framework, we are able to generate the posterior model space measure, $m_{\Delta I}^{p}$, which can be inspected to understand the new inversion method. Such inspections can include simple visual inspection of the posterior model expectation compared to the true field is is inverting for. Both of these can be pushed forward through the sea level operator to obtain their respective sea level fingerprints for comparison.

The covariance between the major sources can be inspected to understand how well the method is able to disambiguate between them.

\subsubsection{Comparison with previous method}
\notebookcite{thomasholland06SimpleInversionipynb2026}

To compare our method's estimate of GMSLC with the altimetry estimate, we can compare the $z$-scores of the two methods' estimates of GMSL change relative to the true GMSL change. We generate a true GMSL value, $\gmslctrue$, and the inversion method's estimate, $\gmslcbay$ by applying the mass balance GMSL operator, $\operatorgmsl{\text{true}}$, to the true ice thickness change field and the posterior model expectation, respectively:

\begin{align}
	\gmslctrue &= \operatorgmsl{\text{true}} \thicknesschange{i}^\dagger, \label{eq:true-gmsl-model} \\
	\gmslcbay &= \operatorgmsl{\text{true}} \thicknesschange{i}^{p}. \label{eq:bay-gmsl-model}
\end{align}


The GMSL altimetry estimate is calculated from the available sea surface height change (SSHC) point measurements, $d_i$, using an area-weighted scheme (\autoref{eq:gmsl_ssh_alt_integral}). Each individual observation is treated as an independent 1D Gaussian measure

\begin{equation}
  d_i \sim \mathcal{N}(d_i, \covariance{n,\,i}),
\end{equation}

where $\covariance{n,\,i}$ represents the noise variance at that specific point. These individual measures are then combined via a cosine-latitude weighted average (with weights $w_i$) to form a single global distribution. Mathematically, this estimate represents the true ice state evaluated through the forward operators plus an aggregated noise term, yielding the following distribution:

\begin{align}
 \gmslcalt &= \operatorgmsl{\text{point}} \circ \operatorsshc \circ \operatorslc \circ \operatorloadice (\thicknesschange{i}^\dagger) + n_{\text{alt}} \\
 &\sim \mathcal{N}\left( \sum_{i=1}^{N} w_i d_i, \sum_{i=1}^{N} w_i^2 \covariance{n,\,i} \right) \label{eq:alt-gmsl-model}
\end{align}

The $z$-score is calculated using the expectation and standard deviations, $\mu$ and $\sigma$ respectively, for both the old altimetry method and the new inversion method, to compare the performance of the two methods in estimating the true GMSL change:

\begin{align}
    z^{bay} &= \frac{\gmslctrue - \mu_{\text{bay}}}{\sigma_{\text{bay}}} \\[7pt]
    z^{alt} &= \frac{\gmslctrue - \mu_{\text{alt}}}{\sigma_{\text{alt}}}
\end{align}

The inversions was run multiple times to generate a suite of $z$-scores for each method, which can be compared to show the improvement of the new method over the old method.

\subsubsection{Sensitivity tests}
\notebookcite{thomasholland07InversionSensitivityipynb2026}

With any inverse setup, it is important to understand how sensitive the results are to the various parameters of the method. This is particularly significant for our method, as it relies on a number of assumptions and choices in the construction of the prior model space measure.

To test the inversion method, the scheme described in \textcite{heathcoteScaleableBayesianFramework2026} for testing the inversion sensitivity to changes in the ice lengthscale, covariance, and expectation bias parameters of the prior.

Initially, an inversion scenario is set up similar to the inversion, with a fixed truth ice state, $m_{\Delta I}^\dagger$. Inversions for the true field were run with priors generated with the systematically varied parameters of interest, and comparison is made between their results and the true field to assess the method's sensitivity to that parameter.

In addition to the above tests from \textcite{heathcoteScaleableBayesianFramework2026}, the sensitivity of the method to the ice thickness prior model to spatial patters or not. True ice states for both spatially uniform ice change and thickness dependent change patters were generated (see \autoref{sec:ice_thickness_priors}). These were both inverted for using the two types of prior model, and then compared.

\subsubsection{Joint inversion using ice, firn, and ocean dynamic changes}
% Paragraph: Describe the joint inversion -- ice + firn + ocean dynamics
% inverted simultaneously. Explain why this matters (correlated signals).

The model space would be a direct sum of the ice thickness change, firn thickness change space, and ocean dynamic topography change:

\begin{equation}
	m = \thicknesschange{i} \oplus \thicknesschange{f} \oplus \thicknesschange{odt}
\end{equation}

With the data space being the direct sum of SSH altimetry data, tide gauge data, and ice altimetry data:

\begin{equation}
	d = d_{\Delta SSH} \oplus d_{\Delta TG} \oplus d_{\Delta ISH}
\end{equation}

It is often the case where multiple model space parameters map to one data space parameter, for example in our case, SSH altimetry points forms the linear equation:

\begin{equation}
	d_{\Delta SSH} = \operatorpointocean( \operatorsshc \circ \operatorslc (\operatorloadice m_{\Delta I} + \operatorloadfirn m_{\Delta F} + \operatorloadwater \odtchange) + \odtchange)
\end{equation}

The forward operator for the joint inversion can be written as a block matrix, with each block being the mapping of the respective model space parameter to the respective data space parameter. This results in the following form:

\begin{align}
\begin{bmatrix}
	\text{SSH Altimetry Points}\\
	\text{Tide Gauge Points}\\
	\text{Ice Altimetry Points}
\end{bmatrix}
&\nonumber\\&=
\begin{bmatrix}
	\operatorpointocean \operatorsshc \operatorslc \operatorloadice & \operatorpointocean \operatorsshc \operatorslc \operatorloadfirn & \operatorpointocean(\operatorsshc \operatorslc \operatorloadwater + 1) \\
	\operatorpointtg \operatorslc \operatorloadice & \operatorpointtg \operatorslc \operatorloadice & \operatorpointtg(\operatorslc \operatorloadwater + 1)\\
	\operatorpointice & \operatorpointice & \mathbf{0}
\end{bmatrix}
\begin{bmatrix}
	\thicknesschange{\text{Ice}}\\
	\thicknesschange{\text{Firn}}\\
	\odtchange
\end{bmatrix}
\end{align}

However, this is computationally inefficient, as it requires six independent calls to the fingerprint numerical solver. To optimise this forward operator to reduce the number of fingerprint calls, we can factorise this matrix into sparse matrices (see \autoref{sec:joint_operator_factorisation}), shown in \autoref{eq:factorised_joint}. This reduced each iteration time, with a speed-up of $3\times$.

\begin{equation}
=
	\begin{bmatrix}
\operatorpointocean \operatorsshc & \operatorpointocean & 0 & 0 \\
\operatorpointtg & \operatorpointtg & 0 & 0 \\
0 & 0 & \operatorpointice & \operatorpointice
\end{bmatrix}
\begin{bmatrix}
\operatorslc & 0 & 0 & 0 \\
0 & \mathbf{1} & 0 & 0 \\
0 & 0 & \mathbf{1} & 0 \\
0 & 0 & 0 & \mathbf{1}
\end{bmatrix}
\begin{bmatrix}
\mathbf{L_I} & \mathbf{L_F} & \mathbf{L_W} \\
0 & 0 & \mathbf{1} \\
\mathbf{1} & 0 & 0 \\
0 & \mathbf{1} & 0
\end{bmatrix}
\begin{bmatrix}
\thicknesschange{\text{Ice}} \\
\thicknesschange{\text{Firn}} \\
\odtchange
\end{bmatrix}\label{eq:factorised_joint}
\end{equation}

% Paragraph: Discuss preconditioning and why it is needed for numerical
% efficiency when the operator is ill-conditioned.


% Paragraph: Describe application to altimetry data and/or DUACS SSH data.

\subsubsection{Sensitivity tests of joint inversions}

TKTK covariance tests

TKTK knock out tests


\section{Results}
\subsection{Error Quantification}
\begin{figure}
\centering
\includegraphics[width=0.45\columnwidth]{01-SSH/load.png}

\includegraphics[width=0.45\columnwidth]{01-SSH/slc.png}
\includegraphics[width=0.45\columnwidth]{01-SSH/ssh.png}
	
	\hfill\includegraphics[width=0.45\columnwidth]{01-SSH/obs.png}\hspace{0.05\columnwidth}
	
	\includegraphics[width=0.8\columnwidth]{01-SSH/global_error_plt.png}
	\caption{A}
	\label{fig:simple_load_to_obs}
\end{figure}

\subsubsection{Error Variations with Altimetry Latitude}

\begin{figure}[!h]
	\includegraphics[width=\columnwidth]{major_source_altimetry_errors_scalar.png}
	\caption{A}\label{fig:simple_major_source}
\end{figure}

\begin{figure}[!h]
	\includegraphics[width=\columnwidth]{figs/double_distribution_66_80}
	\caption{A}\label{fig:gaussian_major_source}
\end{figure}


\subsection{Inversion for Ice Thickness Change from Ocean Altimetry}


\subsection{Joint inversion for Ice and Firn Thickness from Ocean and Ice Altimetry}

\subsection{}

\section{Discussion}
\subsection{Error of using SSHC as a proxy for SLC in GMSL estimation}

The key result from the per-source error characterisation shows that the proxy error differs in both sign and magnitude between ice sources: Greenland melt leads to a systematic underestimation of GMSLC, while West and East Antarctic sources produce a smaller bias of opposite sign. This source dependence means that the SSHC proxy error cannot be corrected by a single scalar factor; the correction would itself depend on the unknown spatial distribution of ice change. This matches the findings of \textcite{lickleyBiasEstimatesGlobal2018, al-attarReciprocitySensitivityKernels2023}.

At the global scale the resulting error is modest, on the order of 1-10\% depending on expected value and uncertainty of the GMSLC; however, regional averages exhibit substantially larger and spatially variable biases, since SSHC and SLC differ by the solid-earth deformation and geoid perturbation terms, which are functions of the ice load which is spatially heterogeneous. This motivates a more principled inference framework in which the full fingerprint physics is incorporated into the estimation procedure.


\subsection{Inversion for ice thickness change from ocean altimetry}

The spatial recovery of the simple inversion is good with GMSLC estimates within one standard deviation
%TC:ignore
(e.g. \autoref{fig:simple_inversion_example}: true GMSLC = -5.47 mm; posterior expectation = -5.50 mm, posterior standard deviation = 0.0456 mm),
%TC:endignore
showing that the Bayesian framework is able to effectively use the ocean altimetry data to constrain the ice thickness change. The sensitivity tests suggest that even with poor calibration of the priors, the ocean altimetry is informative enough to constrain the ice thickness change, which is encouraging for the application of this method to real data where the priors are likely to be imperfect. The repeatability of this across multiple inversion suggests that this is not a feature of a single true data source, but rather a general feature of the method.

The improved $z$-scores compared to the proxy method confirm that the Bayesian inversion is an improvement for estimating GMSLC. The proxy's large $z$-scores reflect systematic uncertainty underestimation: the SSHC proxy cannot represent the full SLC response, whereas the Bayesian framework can, producing well-calibrated posteriors. Well-calibrated uncertainty is practically significant: overly confident GMSLC estimates propagate through to sea level budget closure and coastal hazard projections, where they can mask discrepancies and mislead risk assessments. This was confirmed by the number of true GMSLC values falling within one posterior standard deviation: 60.5\% under the Bayesian method compared to 14.6\% under the proxy method (n=261).

The computational cost of the inference is sensitive to the choice of prior covariance amplitude: conjugate gradient iteration counts ranged from 47 (amplitude multiplier 0.5$\times$) to 610 (multiplier 5.0$\times$), while length-scale and mean-offset variations had comparatively little effect on convergence.

\subsection{Joint inversion using ice, firn, and ocean dynamic changes}

The joint inversion show good recovery of the detail and good sensitivity when using all data types. There are clear trade-offs between the model parameters, as expected, and the method does well to disentangle these variable sources. The lower recovery of firn in some regions is expected due to low sensitivity of mass sensitive data sources to firn changes and could be masked by the more dominant ice changes.

Knock-out tests show that there is information gained from each data type, especially ice altimetry constraining the trade-off between the different densities of firn and ice effectively. SSH altimetry is critical for constraining ocean dynamic topography, which makes intuitive sense as it provides the spatial coverage needed to separate ocean dynamic signals from ice-driven fingerprints.

The loss in ocean dynamic topography detail is acceptable, given the altimetry noise parameter (1\,mm standard deviation) and so only recovering features over this noise level is expected. This is also in part a consequence of jointly inverting for multiple sources, which introduces additional degrees of freedom which lead to some detail loss.


\section{Conclusion}

% For the conclusion, given that the conclusion is currently empty, here is a suggested structure rather than text:
% First, a brief restatement of the problem framing - why existing methods are limited and what this work set out to do differently. One short paragraph, no new information.



% Second, a synthesis of the three main results in order: the error quantification result (with the key numbers - the 5.8-8.2% GIS underestimation, the source-dependence that precludes a scalar correction), the simple inversion result (the 60.5% vs 14.6% within-1-sigma comparison is a strong quantitative claim), and the joint inversion result (what the knockout tests specifically showed about data type contributions).

% Third, a brief paragraph on what the work collectively demonstrates at a higher level - specifically the argument that the Bayesian infinite-dimensional framework gives you things that neither the old proxy method nor the Coulson et al. mascon approach provide: principled uncertainty propagation, spatially-correlated posteriors, and joint inference over competing signals.




This method shows promise, and combined with the underlying work within \textcite{heathcoteScaleableBayesianFramework2026} and the ongoing work within \citetitle{al-attarPySLFPPythonSea2025}, the implementation of this method would be a significant step forward in understanding ice change driven sea level change, and would provide a more principled framework for incorporating the full physics of sea level fingerprints into the estimation process.

However, there are areas of future work that would need to be explored to further develop the method and ensure its applicability to real data. Namely, investigating the sensitivity of the joint inversion to the choice of priors, especially the density choices and spatial fields, and understanding the trade-offs that occur between the thickness changes. This would help to ensure that the method is robust to different assumptions and can be applied to a wide range of scenarios.

Also, there is the possibility of using these methods to recover wider oceanographic details, such as the thermosteric and halosteric effects \autocite{heathcoteScaleableBayesianFramework2026}, which would be useful for understanding the ocean dynamic contributions to sea level change. This would be a powerful tool for understanding the drivers of sea level change, and would provide a more comprehensive picture of the contributions from different sources and could lead to a more complete understanding of the sea level budget.

Of course, implementing the method proposed within this report with real-world data would be the ultimate test of the method, and would be a natural next step for future work. As a proof of concept, a test was run with the \textcite{copernicusclimatechangeserviceGLOBALOCEANGRIDDED2021} dataset, and found that the method can be applied to real data. But without the opportunity to explore the results of this in any detail nor apply the necessary corrections, these would not be presentable. However, the outcome from a subsampled version of the dataset is promising and matches other studies predictions, and suggests that the method can be successfully applied to real data with some further work.


%TC:ignore
\section{Acknowledgements}

I would like to first of all thank my supervisor, David Al-Attar. He has been incredibly helpful and has provided me with much guidance. He has truly gone above and beyond teaching me the maths, far beyond the scope of this project and report so that I had a thorough understanding of the underlying maths. David has been a phenomenal supervisor and provided motivation and support, and has helped me learn a great deal about the world of computational geophysics.

Of course, without David and Daniel Heathcote's work on \citetitle{al-attarPySLFPPythonSea2025} and \citetitle{al-attarPygeoinfPythonLibrary2025}, none of this project would have been possible. Dan has been particularly helpful with oceanographic aspects that I was less familiar with.

I would also like to thank Helen Soane and Cara Hardwick for proofreading the drafts of this report and nodding along to my ramblings on about ice changes and satellite altimetry for the last five months. They have helped enormously with keeping me motivated and have supported me immensely.

Lastly, thanks to the fellow occupants of the Part III Geophysics lab, Peter Jackson, Candela Gil, Tom Devenish-Arzuza, Jiarui Sheng, and Oliver Tang. They have been great company and fun to bounce ideas off and have also listened to my tangents about optimisation and coding best practices and colormaps.

{\let\clearpage\relax
\addcontentsline{toc}{section}{Further Reading}
\section*{Further Reading}}

Readers seeking deeper background on the mathematical and statistical frameworks employed in this report may find the following texts useful: \citetitle{stuartInverseProblemsBayesian2010} \autocite{stuartInverseProblemsBayesian2010} and \textcite{bui-thanhComputationalFrameworkInfiniteDimensional2013} for Bayesian inversion in infinite dimensions; \citetitle{parkerGeophysicalInverseTheory1994} \autocite{parkerGeophysicalInverseTheory1994} for geophysical inverse theory; and \citetitle{rudinFunctionalAnalysis1991} \autocite{rudinFunctionalAnalysis1991} for functional analysis and Hilbert space theory.

\printbibliography[heading=bibintoc]


\appendix
\renewcommand{\thefigure}{A\arabic{figure}}
\setcounter{figure}{0}

\section{Derivations}
\subsection{Sea Level and Sea Surface Height Equations}\label{sec:sea_level_equations}

This derivation is not a new result, but is included here for completeness
and to clarify the assumptions and approximations made in deriving the
equations for sea level change (SLC) and sea surface height change (SSHC)
in response to ice mass redistribution. The derivation follows the framework
laid out in \textcite{al-attarReciprocitySensitivityKernels2023}.

Throughout, we work within a first-order perturbation theory, assuming that
changes in sea level, displacement, and gravitational potential are all small
relative to the background state. We further assume that the shoreline
geometry is fixed, meaning the ocean function, $C$, does not evolve as sea
level changes; this is what renders the fingerprint problem linear and is
appropriate for the small-amplitude changes associated with modern ice melt.
Finally, we assume that the elastic response of the Earth dominates over
viscous relaxation, which is justified for the timescales relevant to
present-day ice mass flux.

Sea level (SL) is defined upon the equipotential surface such that
%
\begin{equation}
    \Phi(\mathbf{x} + \mathrm{SL}\,\hat{\mathbf{r}}) = \Phi_0,
    \label{eq:sl_definition}
\end{equation}
%
where $\Phi$ is the gravitational potential, $\mathbf{x}$ is the position
vector, $\hat{\mathbf{r}}$ is the unit vector in the radial direction, and
$\Phi_0$ is a constant representing the potential at the sea surface. This
means that the sea surface must adapt to changes in the gravitational and
centrifugal potential to maintain an equipotential surface.

When perturbed by an ice thickness change $\Delta T_I$, four things occur
simultaneously:
%
\begin{itemize}
    \item the gravitational potential itself changes, by $\phi$;
    \item the centrifugal potential changes due to perturbations in
          the Earth's rotation vector, by $\psi$;
    \item the solid Earth surface, $\mathbf{x}$, moves radially by
          the displacement vector $\mathbf{u}$; and
    \item the value of the potential at the ocean surface shifts,
          from $\Phi_0$ to $\Phi_0 + \Phi_g$.
\end{itemize}
%
This results in the SL changing, from SL to $\mathrm{SL} + \Delta\mathrm{SL}$,
in the form of
%
\begin{equation}
    (\Phi + \phi + \psi)\!\left(
        \mathbf{x} + \mathbf{u} + (\mathrm{SL} + \Delta\mathrm{SL})\,\hat{\mathbf{r}}
    \right) = \Phi_0 + \Phi_g,
    \label{eq:sl_perturbed}
\end{equation}
%
to which we apply a first-order Taylor expansion of the left-hand side to get
%
\begin{equation}
    \Phi(\mathbf{x} + \mathrm{SL}\,\hat{\mathbf{r}})
    + \phi + \psi
    + \nabla\Phi \cdot \left(\mathbf{u} + \Delta\mathrm{SL}\,\hat{\mathbf{r}}\right)
    = \Phi_0 + \Phi_g.
    \label{eq:sl_taylor}
\end{equation}
%
This equation can be simplified using:
%
\begin{itemize}
    \item using the relation in \autoref{eq:sl_definition} we can cancel the
          terms involving $\Phi$ and $\Phi_0$;
    \item $\nabla\Phi$ represents the gravitational acceleration vector with
          its radial magnitude as $g$, therefore
          $\nabla\Phi \cdot \hat{\mathbf{r}} = g$; and
    \item the term $\mathbf{u} \cdot \nabla\Phi$ represents the potential
          change due to the vertical motion of the solid surface;
\end{itemize}
%
to give
%
\begin{equation}
    \phi + \psi + \mathbf{u} \cdot \nabla\Phi + g\,\Delta\mathrm{SL} = \Phi_g,
    \label{eq:sl_simplified}
\end{equation}
%
which we can rearrange to get the change in sea level
%
\begin{equation}
    \Delta\mathrm{SL} = -\frac{1}{g}
        \left(\mathbf{u} \cdot \nabla\Phi + \phi + \psi\right)
        + \frac{\Phi_g}{g}.
    \label{eq:slc}
\end{equation}

When calculating SSH, we need to account for the fact that SSH is the
height of the sea surface above a reference ellipsoid, meaning we must
add back the radial displacement of the solid surface. We additionally
choose to remove the centrifugal contribution $\psi$ from the SSH change,
following \textcite{al-attarReciprocitySensitivityKernels2023}; this is a convention motivated by the fact
that satellite altimetry measures height above a reference ellipsoid fixed
to the mean rotational state, so that isolating $\phi$ gives a cleaner
measure of the mass redistribution signal. Whether and how this should be
done is a subtle point discussed in detail in \textcite{tamisieaOngoingGlacialIsostatic2011} and
\textcite{lickleyBiasEstimatesGlobal2018}. Under this convention, SSHC is related to
SLC by
%
\begin{equation}
    \Delta\mathrm{SSH} = \Delta\mathrm{SL}
        + \hat{\mathbf{r}} \cdot \mathbf{u}
        + \frac{\psi}{g}.
    \label{eq:ssh_from_sl}
\end{equation}
%
Since $\nabla\Phi$ points approximately radially with magnitude $g$ at the
Earth's surface, we make the approximation
$\mathbf{u} \cdot \nabla\Phi \approx g\,(\hat{\mathbf{r}} \cdot \mathbf{u})$,
which neglects the contribution of lateral surface displacement to the
potential change. Substituting \autoref{eq:slc} into \autoref{eq:ssh_from_sl}:
%
\begin{equation}
    \Delta\mathrm{SSH} =
    \underbrace{
        -\frac{1}{g}\!\left(
            g\,(\hat{\mathbf{r}} \cdot \mathbf{u}) + \phi + \psi
        \right)
        + \frac{\Phi_g}{g}
    }_{\Delta\mathrm{SL}}
    + \hat{\mathbf{r}} \cdot \mathbf{u}
    + \frac{\psi}{g},
    \label{eq:ssh_expanded}
\end{equation}
%
in which the $\hat{\mathbf{r}} \cdot \mathbf{u}$ terms cancel and the
$\psi/g$ terms cancel, which simplifies to
%
\begin{equation}
    \Delta\mathrm{SSH} = \frac{\Phi_g}{g} - \frac{\phi}{g}.
    \label{eq:sshc}
\end{equation}
%
This result is physically intuitive: SSH change has spatial structure
through $\phi$, the perturbation to the gravitational potential due to
ice mass redistribution, but unlike SL it has no dependence on solid
surface displacement $\hat{\mathbf{r}} \cdot \mathbf{u}$ or centrifugal
potential perturbation $\psi$.

\subsection{Derivation of the relationship between GMSL approximation and SSHC}\label{sec:point_gmsl_derivation}

The motivation of this section is to derive the mathematical relationship between the GMSL approximation and the SSHC values across the altimetry domain, which is used in the traditional method of estimating GMSL change from ocean altimetry data. By defining weights for each datapoint based on their latitude, we can express the GMSL approximation as a weighted sum of the SSHC values using an averaging operator built into \citetitle{al-attarPygeoinfPythonLibrary2025}.

To find the GMSL approximation with regards to SSHC, from \autoref{eq:gmsl_ssh_integral} and \autoref{eq:gmsl_ssh_alt_integral}, we start by considering the integral of the SSHC across a section of the Earth's surface, which we denote as $\Omega$ which has an area of $|\Omega|$:


\begin{equation}
    \gmslc = \dfrac{1}{|\Omega|}\int_\Omega \sshc(\phi, \lambda) \, dA \label{eq:gmsl_sshc_integral}
\end{equation}

Where $\sshc(\phi, \lambda)$ is the sea surface height change at a given longitude $\phi$ and latidue $\lambda$. With a spherical radius of $R$, the area, $dA$, is non-uniform in size and shrinks towards the poles, and can be expressed as:

\begin{align}
    dA &= (R\cos \lambda d\phi)(R d\lambda) \\
    &= R^2 \cos \lambda \, d\phi \, d\lambda
\end{align}

When we take discrete samples across a uniform grid, with spacings of $\Delta \phi$ and $\Delta \lambda$, we convert \autoref{eq:gmsl_sshc_integral} into a double summation over N points:

\begin{equation}
  \gmslc \approx \dfrac{\sum^N_{i=1}\sshc(\mathbf{x}_i)\cdot(R^2\cos \lambda_i \Delta \phi \Delta \lambda)}{\sum^N_{j=1}(R^2 \cos \lambda_j\Delta \phi\Delta \lambda)}
\end{equation}

Which can simplify to:

\begin{equation}
  \gmslc \approx \sum^N_{i=1}\sshc(\mathbf{x}_i)\cdot\dfrac{\cos \lambda_i}{\sum^N_{j=1}\cos \lambda_j}
\end{equation}

Which for ease within an averaging quadrature scheme, we can define the weights for each width, $w_i$, as a ratio of it's local cosine width to the sum of all the cosine widths of the avaliable data points:

\begin{equation}
    w_i = \dfrac{\cos \lambda_i}{\sum^N_{j=1}\cos \lambda_j}
\end{equation}

Which then leads to the final expression for the GMSL approximation as a weighted sum of the SSHC values across the altimetry domain:

\begin{equation}
    \gmslc \approx \sum^N_{i=1} w_i \sshc(\mathbf{x}_i)
\end{equation}


\subsection{Joint operator factorisation}\label{sec:joint_operator_factorisation}

We have the forward operator matrix:

\begin{equation}
    O = \begin{bmatrix}
        \operatorpointocean \operatorsshc \operatorslc \operatorloadice & \operatorpointocean \operatorsshc \operatorslc \operatorloadfirn & \operatorpointocean(\operatorsshc \operatorslc \operatorloadwater + 1) \\
        \operatorpointtg \operatorslc \operatorloadice & \operatorpointtg \operatorslc \operatorloadfirn & \operatorpointtg(\operatorslc \operatorloadwater + 1) \\
        \operatorpointice & \operatorpointice & 0 \\
        \operatorgrace \operatorslc \operatorloadice & \operatorgrace \operatorslc \operatorloadfirn & \operatorgrace \operatorslc \operatorloadwater
    \end{bmatrix}
\end{equation}

Which maps the model space vector $\mathbf{x} = (\thicknesschangeice,\, \thicknesschangefirn,\, \odtchange)$ to the data space vector $d = (d_{\Delta SSH},\, d_{\Delta TG},\, d_{\Delta ISH},\, d_{\text{GRACE}})$.

This is an inefficient operator, as it requires six independent calls to the fingerprint numerical solver $\operatorslc$, which is the most computationally expensive part of the forward model, due to its numerical implementation.

We can factorise this matrix into three sparse matrices $A$, $B$, and $C$ such that $O = ABC$, where $B$ contains only one call to $\operatorslc$, thus reducing the number of fingerprint calls from six to one. The derivation of this factorisation proceeds in three steps:

\textbf{Step 1: Isolate the fingerprint structure}

Notice that $\operatorslc$ appears applied to $\operatorloadice$, $\operatorloadfirn$, $\operatorloadwater$
in rows 1, 2, and 4, but never in row 3. This suggests introducing an intermediate quantity
$\mathbf{u} = \operatorslc\mathbf{L}\,\mathbf{x}$ where $\mathbf{L}$ collects the load operators,
so that $\operatorslc$ is called only once. Define

\begin{equation}
    C = \begin{bmatrix}
        \operatorloadice & \operatorloadfirn & \operatorloadwater \\
        0 & 0 & 1 \\
        1 & 0 & 0 \\
        0 & 1 & 0
    \end{bmatrix}
\end{equation}

where the bottom three rows simply route the state variables
$(\thicknesschangeice,\, \thicknesschangefirn,\, \odtchange)$ to slots that will be
used by the outer operators. The action of $C$ on $\mathbf{x}$ gives the intermediate vector

\begin{equation}
    C\mathbf{x} = \begin{pmatrix}
        \operatorloadice\,\thicknesschangeice + \operatorloadfirn\,\thicknesschangefirn + \operatorloadwater\,\odtchange \\
        \odtchange \\
        \thicknesschangeice \\
        \thicknesschangefirn
    \end{pmatrix}
\end{equation}

\textbf{Step 2: Apply the fingerprint operator selectively}

We want $\operatorslc$ applied to only the first component of $C\mathbf{x}$, leaving the rest
unchanged. This is exactly what a block-diagonal operator achieves:

\begin{equation}
    B = \begin{bmatrix}
        \operatorslc & 0 & 0 & 0 \\
        0 & 1 & 0 & 0 \\
        0 & 0 & 1 & 0 \\
        0 & 0 & 0 & 1
    \end{bmatrix}
\end{equation}

so that

\begin{equation}
    BC\mathbf{x} = \begin{pmatrix}
        \operatorslc(\operatorloadice\,\thicknesschangeice + \operatorloadfirn\,\thicknesschangefirn + \operatorloadwater\,\odtchange) \\
        \odtchange \\
        \thicknesschangeice \\
        \thicknesschangefirn
    \end{pmatrix}
\end{equation}

Linearity of $\operatorslc$ then distributes the first component into
$\operatorslc\operatorloadice\,\thicknesschangeice + \operatorslc\operatorloadfirn\,\thicknesschangefirn
+ \operatorslc\operatorloadwater\,\odtchange$, which is precisely the fingerprint contribution
needed in rows 1, 2, and 4 of $M$.

\textbf{Step 3: Recover the original rows by projection}

It remains to map the four-component vector $BC\mathbf{x}$ back to the four observation types.
Reading off what each row of $M$ requires:
\begin{itemize}
    \item Row 1 needs $\operatorpointocean \operatorsshc$ acting on the fingerprint result, plus $\operatorpointocean$ acting on $\odtchange$ directly.
    \item Row 2 needs $\operatorpointtg$ acting on the fingerprint result, plus $\operatorpointtg$ acting on $\odtchange$ directly.
    \item Row 3 needs $\operatorpointice$ acting on $\thicknesschangeice$ and $\thicknesschangefirn$, with no fingerprint contribution.
    \item Row 4 needs $\operatorgrace$ acting on the fingerprint result only. Unlike rows 1 and 2, there is no direct additive contribution from $\odtchange$, since $\operatorgrace$ extracts the gravitational potential perturbation $\phi$ from the fingerprint response space, which is a globally defined field insensitive to the geometric height redistribution of ODT.
\end{itemize}

This gives

\begin{equation}
    A = \begin{bmatrix}
        \operatorpointocean \operatorsshc & \operatorpointocean & 0 & 0 \\
        \operatorpointtg & \operatorpointtg & 0 & 0 \\
        0 & 0 & \operatorpointice & \operatorpointice \\
        \operatorgrace & 0 & 0 & 0
    \end{bmatrix}
\end{equation}

\textbf{Conclusion}

Combining the three steps, $O = ABC$, where each factor has a clear role: $C$ computes the total load and routes state variables, $B$ applies $\operatorslc$ exactly once to the combined load, and $A$ projects the result onto the four observation spaces.


\section{Supplementary figures}

\begin{figure}[h]
  \begin{center}
  \includegraphics[width=0.9\columnwidth]{7-0_baseline_inversion.pdf}
  \end{center}

	\caption{}
	\label{fig:sensitivity-true-inversion}

\end{figure}



\begin{figure}[h]
  \begin{center}
  \includegraphics[width=0.9\columnwidth]{7-1_length_scale.pdf}
  \end{center}

	\caption{}
	\label{fig:sensitivity-lengthscale}

\end{figure}


\begin{figure}[h]
  \begin{center}
  \includegraphics[width=0.9\columnwidth]{7-2_mean_translation.pdf}
  \end{center}

	\caption{}
	\label{fig:sensitivity-offset}

\end{figure}


\begin{figure}[h]
  \begin{center}
  \includegraphics[width=0.9\columnwidth]{7-3_covariance_scaling.pdf}
  \end{center}

	\caption{}
	\label{fig:sensitivity-variance}

\end{figure}


\section{Mathematical notation}
% ============================================================
%  Mathematical Notation
% ============================================================

\subsection*{Physical fields and scalar quantities}

\begin{tabular}{@{}lp{10cm}@{}}
    \toprule
    \textbf{Symbol} & \textbf{Description} \\
    \midrule
    $\mathrm{SL}(\mathbf{x})$   & Sea level (relative sea level) at position $\mathbf{x}$ \\
    $\mathrm{SSH}(\mathbf{x})$  & Sea surface height at position $\mathbf{x}$ \\
    $\slc$                       & Sea level change \\
    $\sshc$                      & Sea surface height change \\
    $\gmsl$                      & Global mean sea level \\
    $\gmslc$                     & Global mean sea level change \\
    $\gmslctrue$                 & True (mass-balance) GMSL change \\
    $\gmslcbay$                  & Bayesian inversion estimate of GMSL change \\
    $\gmslcalt$                  & Altimetry-based estimate of GMSL change \\
    $\odtchange$                 & Ocean dynamic topography change \\
    $\mathrm{C}(\mathbf{x})$    & Ocean function; 1 on ocean, 0 on land \\
    $\mathrm{I}(\mathbf{x})$    & Ice thickness at position $\mathbf{x}$ \\
    \bottomrule
\end{tabular}

\vspace{1em}

\subsection*{Model space variables}

\begin{tabular}{@{}lp{10cm}@{}}
    \toprule
    \textbf{Symbol} & \textbf{Description} \\
    \midrule
    $\thicknesschangeice$              & Ice thickness change field \\
    $\thicknesschange{f}$              & Firn thickness change field \\
    $\thicknesschange{odt}$            & Ocean dynamic topography thickness change field \\
    $\thicknesschangeice$            & Ice thickness change (block-matrix notation) \\
    $\thicknesschangefirn$           & Firn thickness change (block-matrix notation) \\
    $\masschange{i}$                   & Ice mass change \\
    $m$                                & Generic model space element \\
    $\bar{m}$                          & Prior mean of model space element \\
    $m_p$                              & Posterior model space element \\
    $m_{\Delta I}^{\dagger}$           & True (synthetic) ice thickness field \\
    $m_{\Delta I}^{p}$                 & Posterior ice thickness field \\
    $\mathbf{x}$                       & State vector; also used as a position on the sphere \\
    \bottomrule
\end{tabular}

\vspace{1em}

\subsection*{Densities, areas, and physical constants}

\begin{tabular}{@{}lp{10cm}@{}}
    \toprule
    \textbf{Symbol} & \textbf{Description} \\
    \midrule
    $\rho_w$            & Density of water \\
    $\rho_i$            & Density of ice \\
    $\rho_{\text{firn}}$& Density of firn \\
    $A$                 & Total ocean surface area \\
    $A_{\text{ocean}}$  & Total ocean surface area (mass-balance form) \\
    $R$                 & Radius of the Earth \\
    $r$                 & Radial coordinate \\
    $\hat{r}$           & Radial unit vector \\
    $\mathbf{u}$        & Solid-earth displacement vector \\
    \bottomrule
\end{tabular}

\vspace{1em}

\subsection*{Domains and geometry}

\begin{tabular}{@{}lp{10cm}@{}}
    \toprule
    \textbf{Symbol} & \textbf{Description} \\
    \midrule
    $\mathbb{S}^2$              & Unit two-sphere (used as a model for the Earth's surface) \\
    $\Omega$                    & Ocean domain (full altimetry-observable ocean) \\
    $\Omega_{\altimetry}$       & Ocean domain within satellite altimetry coverage \\
    $\Omega_I$                  & Ice field domain \\
    $|\Omega|$                  & Surface area of domain $\Omega$ \\
    $\phi$                      & Longitude \\
    $\lambda$                   & Latitude \\
    $dA$                        & Surface area element \\
    $dS$                        & Surface area element (spherical) \\
    $\Delta\phi,\,\Delta\lambda$& Grid spacing in longitude and latitude \\
    $N$                         & Number of discrete sampling points \\
    $w_i$                       & Area-weighting coefficient for point $i$; $w_i = \cos\lambda_i / \sum_j \cos\lambda_j$ \\
    \bottomrule
\end{tabular}

\vspace{1em}

\subsection*{Function spaces}

\begin{tabular}{@{}lp{10cm}@{}}
    \toprule
    \textbf{Symbol} & \textbf{Description} \\
    \midrule
    $H^s(\mathbb{S}^2)$     & Sobolev space of index $s$ on the sphere \\
    $s$                      & Sobolev smoothness index \\
    $h, k$                   & Elements of $H^s(\mathbb{S}^2)$ \\
    $X,\, Y$                 & Generic Hilbert spaces \\
    $\braket{\cdot,\cdot}_Y$ & Inner product on space $Y$ \\
    $\phi$                   & Continuous linear functional $\phi: H^s(\mathbb{S}^2) \to \mathbb{R}$ \\
    \bottomrule
\end{tabular}

\vspace{1em}

\subsection*{Operators}

\begin{tabular}{@{}lp{10cm}@{}}
    \toprule
    \textbf{Symbol} & \textbf{Description} \\
    \midrule
    $\operatorslc$                          & Sea level fingerprint operator; maps load change to sea level change \\
    $\operatorsshc$                         & Sea surface height change operator; maps load change to SSHC \\
    $\operatorloadice$                      & Ice load operator; point-wise multiplication by $\rho_i$ \\
    $\operatorloadfirn$                     & Firn load operator; point-wise multiplication by $\rho_{\text{firn}}$ \\
    $\operatorloadwater$                    & Water load operator \\
    $\operatorgmsl{\cdot}$                  & GMSL averaging operator \\
    $\operatorgmsl{\text{true}}$            & Mass-balance GMSL operator (integral over ice domain) \\
    $\operatorgmsl{\sshc}$                  & Full-ocean SSHC averaging operator \\
    $\operatorgmsl{\altimetry}$             & Altimetry-domain SSHC averaging operator \\
    $\operatorgmsl{\text{point}}$           & Area-weighted point GMSL estimator \\
    $\operatoraltimetry{\cdot}$             & General altimetry point-sampling operator \\
    $\operatorpointocean$                   & SSH ocean altimetry point-sampling operator \\
    $\operatorpointtg$                      & Tide gauge point-sampling operator \\
    $\operatorpointice$                     & Ice altimetry point-sampling operator \\
    $A$                                     & Generic forward operator; $A: X \to Y$ \\
    $A^*$                                   & Adjoint of operator $A$; $A^*: Y \to X$ \\
    $\alpha(\mathbf{x})$                    & Spatial ice/firn partitioning function (generalised logistic curve) \\
    $W$                                     & Diagonal weighting operator; $W = \operatorname{diag}(\alpha(T_i))$ \\
    \bottomrule
\end{tabular}

\vspace{1em}

\subsection*{Probabilistic quantities}

\begin{tabular}{@{}lp{10cm}@{}}
    \toprule
    \textbf{Symbol} & \textbf{Description} \\
    \midrule
    $\mu$                           & Gaussian measure; also used as prior measure \\
    $\mathcal{N}(m, \mathcal{C})$   & Gaussian measure with mean $m$ and covariance operator $\mathcal{C}$ \\
    $\mathcal{C}$                   & Generic covariance operator \\
    $\covariance{0}$                & Base (heat-kernel) covariance operator \\
    $\covariance{m}$                & Prior model covariance operator \\
    $\covariance{n}$                & Noise / observation covariance operator \\
    $\covariance{\thicknesschangeice}$ & Covariance operator for ice thickness change \\
    $\covariance{\thicknesschange{f}}$ & Covariance operator for firn thickness change \\
    $d$                             & Data space vector \\
    $d_{\text{obs}}$                & Observed data vector \\
    $d_{\Delta\mathrm{SSH}}$        & SSH altimetry data vector \\
    $d_{\Delta\mathrm{TG}}$         & Tide gauge data vector \\
    $d_{\Delta\mathrm{ISH}}$        & Ice surface height altimetry data vector \\
    $n$                             & Observation noise vector; $n \sim \mathcal{N}(0, \covariance{n})$ \\
    $K$                             & Kalman gain matrix \\
    $z$                             & $z$-score of the true value under a predicted distribution \\
    $\sigma_{\text{target}}$        & Target standard deviation of GMSL contribution \\
    $\sigma_{\text{gmsl}}$          & Implied GMSL standard deviation of a thickness change prior \\
    $p(\cdot|\cdot)$                & Conditional probability density \\
    \bottomrule
\end{tabular}

\vspace{1em}

\subsection*{Block-matrix and joint inversion}

\begin{tabular}{@{}lp{10cm}@{}}
    \toprule
    \textbf{Symbol} & \textbf{Description} \\
    \midrule
    $\oplus$        & Direct sum of model or data spaces \\
    $M$             & Full joint forward operator matrix \\
    $A,\,B,\,C$     & Sparse factor matrices in the factorisation $M = ABC$ \\
    $\mathbf{L}$    & Combined load operator matrix \\
    \bottomrule
\end{tabular}

\vspace{1em}

\subsection*{Ice/firn partitioning function parameters}

\begin{tabular}{@{}lp{10cm}@{}}
    \toprule
    \textbf{Symbol} & \textbf{Description} \\
    \midrule
    $x$     & Standardised ice thickness value; $x \in [0,1]$ \\
    $a$     & Lower asymptote of the generalised logistic curve \\
    $k$     & Upper asymptote of the generalised logistic curve \\
    $b$     & Steepness parameter of the generalised logistic curve \\
    $m$     & Inflection threshold of the generalised logistic curve \\
    $\nu$   & Asymmetry parameter of the generalised logistic curve \\
    \bottomrule
\end{tabular}



%TC:endignore
\end{document}
